\section{从实模式走向 64 位内核}

\subsection{转换不是一条指令}

进入长模式至少需要验证 CPU 能力、建立 GDT、打开 A20、建立页表、设置
CR4.PAE、写 EFER.LME、设置 CR0.PG 和 CR0.PE，再通过远跳转装载 64 位代码
段。顺序错误会导致通用保护异常或三重故障。

\begin{center}
\begin{tikzpicture}[os diagram]
  \node[os state] (real) {16 位实模式};
  \node[os block,right=of real] (protected) {32 位保护模式};
  \node[os block,right=of protected] (paging) {PAE 四级页表};
  \node[os state,right=of paging] (long) {64 位长模式};
  \draw[os arrow] (real) -- (protected);
  \draw[os arrow] (protected) -- (paging);
  \draw[os arrow] (paging) -- (long);
\end{tikzpicture}
\end{center}

\subsection{A20 与地址环绕}

早期 PC 为兼容 8086 曾让地址第 20 位可屏蔽。若 A20 未开启，1 MiB 以上地址
可能回绕。项目必须自行启用并验证 A20，不能因为 QEMU 默认环境“似乎可用”
就删除这项契约。

\subsection{Stage 1 与磁盘}

v0.2 将通过 IDE PIO 读取自研 Stage 1。驱动需要选择通道、设备和 LBA，
发送命令，轮询 BSY/DRQ/ERR，并按 16 位数据端口读取扇区。每次传输都要有
超时、扇区范围和错误状态验证。

\subsection{ELF64 装载}

ELF 头不是值得盲目信任的结构体。装载器验证 magic、类别、端序、机器类型、
程序头表范围和每个 PT\_LOAD 段；检查文件范围不越界、
\code{filesz <= memsz}、目标区间不重叠保留内存；复制文件字节并清零 BSS，
最后构造 BootInfo 交给内核入口。

\subsection{实模式与保护模式的历史接力}

8086 分段把 16 位段值左移四位再加 16 位偏移，从有限寄存器构造 20 位地址。
这种设计让代码、数据和栈可以分别移动，却没有可靠的权限隔离。80286 引入保护
模式，段寄存器不再直接贡献地址，而是选择 GDT/LDT 中的描述符；描述符包含
基址、界限、类型和特权级。80386 又把偏移和通用寄存器扩展到 32 位，并加入
分页。

现代 64 位内核通常使用基址为零、覆盖完整空间的平坦段，把地址隔离交给页表。
但保护模式的描述符、特权级和门机制仍是进入长模式及建立用户态边界的一部分。
x86-64 因此不是抛弃旧模式的新体系，而是在兼容链上增加的新执行状态。

\subsection{GDT 为模式切换提供代码属性}

全局描述符表由 GDTR 指向。每个段描述符编码基址、界限、可执行/可写属性、
当前特权级和默认操作数宽度。软件先在内存中构造 GDT，用 \code{lgdt} 装载
表位置，再通过远跳转同时更新 CS 选择子与流水线中的代码属性。

只写 CR0.PE 不会自动刷新 CS 的隐藏缓存。远跳转是状态转换的一部分：它让
后续字节按新代码段的规则译码。数据段寄存器随后也要装入有效选择子，不能继续
假定复位值适用于保护模式。

\subsection{A20 兼容门保留了 8086 环绕行为}

8086 的 20 位地址加法会在 1 MiB 处自然回绕。80286 可以产生更高地址后，
部分旧软件仍依赖这一现象，IBM PC/AT 因而增加了可屏蔽地址线 20 的兼容门。
现代硬件通常默认打开 A20，模拟器也可能表现得始终可用，但启动代码应执行并
验证自己的平台契约。

验证不能只写控制位。固件要选择两个只在 A20 开启时不同的安全地址，保存原值、
写入区分模式、比较结果并恢复，且不能覆盖自身、栈或设备内存。

\subsection{四级页表把地址翻译分解为索引}

经典 4 KiB 页面下，48 位规范虚拟地址被拆成四个 9 位索引和 12 位页内偏移：
\[
  9\;(\mathrm{PML4})+9\;(\mathrm{PDPT})+9\;(\mathrm{PD})
  +9\;(\mathrm{PT})+12\;(\mathrm{offset})=48.
\]
每张表占 4 KiB，含 512 个 64 位表项。CR3 保存顶级表物理地址，处理器逐级
读取 present、writable、user 和 execute-disable 等属性，最终得到物理页。

\begin{pdfdiagram}
  \node[os state] (virtual) {48 位规范\\虚拟地址};
  \node[os block,right=of virtual] (pml4) {PML4\\索引 47:39};
  \node[os block,right=of pml4] (pdpt) {PDPT\\索引 38:30};
  \node[os block,below=of pdpt] (pd) {PD\\索引 29:21};
  \node[os block,left=of pd] (pt) {PT\\索引 20:12};
  \node[os hardware,left=of pt] (page) {物理页\\加偏移 11:0};
  \draw[os arrow] (virtual) -- (pml4);
  \draw[os arrow] (pml4) -- (pdpt);
  \draw[os arrow] (pdpt) -- (pd);
  \draw[os arrow] (pd) -- (pt);
  \draw[os arrow] (pt) -- (page);
\end{pdfdiagram}
\diagramnote{页表既做地址翻译，也在每一级累计权限；上级禁止的权限不能由下级重新开放。}

最初页表只映射切换代码、栈、BootInfo 和内核装载区。映射过少会在启用分页
后的下一次取指或栈访问触发页故障；映射全部内存虽容易启动，却会掩盖权限边界。
早期可先使用 2 MiB 大页减少表数量，内核内存管理建立后再细化权限。

\subsection{控制位的依赖顺序}

长模式能力先由 CPUID 验证。软件建立页表并把物理地址写入 CR3，设置 CR4.PAE，
再在 EFER 中设置 LME。随后开启 CR0.PE 与 CR0.PG，处理器进入长模式激活
状态；最后通过指向 L 位代码段的远跳转进入 64 位子模式。

这些位不是相互独立的开关。没有 PAE 页表就开启分页会使地址解释错误；没有
有效 GDT 就远跳转会产生 \#GP；异常处理尚未建立时，连续异常可能升级为双重
故障和三重故障，平台最终复位。每设置一个控制位，代码都要确保下一条取指和
数据访问仍有合法映射。

\subsection{磁盘扇区是装载器的第一个不可信输入}

IDE PIO 通过端口选择设备与 LBA，写入读取命令，然后等待 BSY 清除和 DRQ
置位。状态寄存器中的 ERR 或 DF 必须立即转入失败路径；每次等待都有预算。
数据端口宽度为 16 位，一个 512 字节扇区需要读取 256 个字。

即使磁盘镜像由项目生成，装载器也按不可信输入验证。扇区数量乘法可能溢出，
LBA 可能越过镜像，声明长度可能超过目标缓冲区。把这些错误在固件层拒绝，
能够避免损坏页表、栈和正在执行的代码。

\subsection{ELF 将文件布局与内存布局分离}

ELF 的节表主要服务链接与调试，运行时装载器依据程序头。每个
\code{PT_LOAD} 项分别给出文件偏移、虚拟地址、文件字节数、内存字节数、
对齐和权限。BSS 在文件中不保存全零内容，所以
\code{p_memsz} 可以大于 \code{p_filesz}；装载器复制已有字节，再把余下
范围清零。

所有加法和乘法在使用前做溢出检查，区间采用半开形式并检查相互重叠。入口地址
必须落入可执行装载段，目标区不能覆盖固件、页表、栈和 BootInfo。只有整个
ELF 验证通过后才开始不可逆复制，避免“验证到一半已经破坏系统”的状态。

\subsection{BootInfo 固定固件与内核的交接}

固件知道内存布局、磁盘来源、内核段位置和早期日志设备；内核不应重新猜测这些
事实。BootInfo 用带版本与长度的显式宽度字段交付内存区间、保留区、启动设备
和日志配置。指针字段必须说明是物理地址还是当前页表下的虚拟地址。

交接文档同时规定寄存器、栈对齐、中断状态、页表和对象生命周期。版本不兼容
时内核应在最早日志通道报告并停止，而不是按错误结构继续解析。
